% 04_parameter_design.tex — 第3节 参数设计
\section{参数设计}
\label{sec:parameters}

\subsection{熵权重计算}

各约束严重度的权重$\omega_a, \omega_f, \omega_v$采用熵权法自适应确定，使变异性大的约束获得更高权重：

步骤1：对每种故障场景下的约束值矩阵$\mathbf{F} \in \mathbb{R}^{N \times 3}$进行归一化：
\begin{equation}
    \tilde{f}_{ij} = \frac{f_{ij} - \min_j f_{ij}}{\max_j f_{ij} - \min_j f_{ij}}
    \label{eq:normalize}
\end{equation}

步骤2：计算信息熵$E_j$和差异系数$d_j$：
\begin{equation}
    E_j = -\frac{1}{\ln N} \sum_{i=1}^N p_{ij} \ln p_{ij}, \quad d_j = 1 - E_j
    \label{eq:entropy}
\end{equation}
其中$p_{ij} = \tilde{f}_{ij} / \sum_i \tilde{f}_{ij}$。

步骤3：归一化权重：
\begin{equation}
    \omega_j = \frac{d_j}{\sum_{j=1}^3 d_j}
    \label{eq:entropy_weight}
\end{equation}

熵权法使得在低新能源渗透率下（功角约束变异大），$\omega_a$较高；在高新能源渗透率下（频率约束激活），$\omega_f$增大，实现权重的自适应调整。

\subsection{GP核函数与超参数选择}

Matérn 5/2核的超参数$(\sigma_f^2, \ell, \sigma_n^2)$通过最大化对数边际似然优化：
\begin{equation}
    \log p(\mathbf{y} | \mathbf{X}, \boldsymbol{\theta}) = -\frac{1}{2}\mathbf{y}^T K_{\boldsymbol{\theta}}^{-1}\mathbf{y} - \frac{1}{2}\log|K_{\boldsymbol{\theta}}| - \frac{N}{2}\log 2\pi
    \label{eq:marginal_likelihood}
\end{equation}
采用5次随机重启避免局部最优，初始噪声$\sigma_n^2 = 10^{-4}$。长度尺度$\ell$初始化为各维取值范围的0.5倍。

\subsection{AIA收缩因子设计}

分离超平面添加后，对其法向量方向施加收缩$\epsilon$：
\begin{equation}
    b_i' = b_i - \epsilon \|\bA_i\|
    \label{eq:shrinkage}
\end{equation}

$\epsilon$的取值需平衡保守性与实用性。本文取$\epsilon = 0.01$，约为参数空间范围的1\%，确保边界不排除已知安全点的同时提供足够的安全裕度。

\subsection{闭环迭代收敛准则}

闭环迭代终止条件：
\begin{equation}
    \frac{|V^{(r)} - V^{(r-1)}|}{V^{(r-1)}} < \eta
    \label{eq:terminate}
\end{equation}
其中$\eta = 0.05$为相对体积变化容差。当连续两轮迭代的安全域体积增长不超过5\%时，认为边界已收敛。
